\section{Structure and certificates}\label{sec:theory}

\subsection{Policy inclusion does not imply diminishing returns}
\begin{proposition}[Calendar inclusion]\label{prop:inclusion}
For fixed $A$ and $R$, $S\subseteq U$ implies $V(S,R)\le V(U,R)$.
Consequently $W_K(R)$ is nondecreasing, with $W_0=V_0$ and $W_{T-1}=V_F$.
\end{proposition}
\begin{proof}
An $S$-policy can be implemented under $U$ by choosing the corresponding
parts of each original block and ignoring observations at additional update
epochs. The player remembers the original block or the privately sampled
pure policy. This embedding is outcome-equivalent against every policy in
the same $\QQ_R$. Extending the embedding to mixtures and maximizing over
the enlarged feasible set proves the value inequality. The budget and
endpoint statements follow immediately.
\end{proof}
This is a finite policy-inclusion argument, consistent with established
zero-sum information-value results \citep{hogeboom2020}. We use it as a
correctness condition, not as an independent novelty claim. It does not
compare different physical speeds or different opponent classes.

\begin{proposition}[Complementary update opportunities]\label{prop:counterexample}
Even with binary actions, a committed opponent and a nonnegative bounded
payoff, $S\mapsto V(S,C)$ need not be submodular.
\end{proposition}
\begin{proof}
Take $T=3$ and payoff
$A=\mathbf1\{b_1=r_0\}\mathbf1\{b_2=r_1\}$.
The values for $S=\varnothing,\{1\},\{2\},\{1,2\}$ are
$1/4,1/2,1/2,1$, respectively. Without updates, independent uniform Blue
guesses guarantee $1/4$, and independent uniform Red bits bound the value
by $1/4$. With only update 1, Blue matches the observed first bit and guesses
the second; with only update 2 it guesses the first and matches the second.
Uniform unknown bits bound both values by $1/2$. Both updates permit both
matches. Thus the marginal value of update 2 rises from $1/4$ to $1/2$
when update 1 is already present, violating submodularity.
\end{proof}
This example concerns the general finite history-payoff class, not every
formation instance. It rules out a generic diminishing-returns justification
for greedy calendar selection. The informational-substitutes literature
likewise requires substantive conditions on the decision problem
\citep{chen2017}. A sufficient special case is a sum of independent reset
games with product feasible sets, independent information and no shared
opponent resource: product equilibrium strategies give matching lower and
upper guarantees for the sum, so each selected update contributes a fixed
premium. That special case is modular. The formation history in
\eqref{eq:formation} generally prevents such a reset decomposition.

\subsection{Where the game can be decomposed}
\begin{theorem}[Common-update decomposition]\label{thm:decomposition}
At a common update $c\in S\cap R$, each complete executed prefix defines a
proper subgame. Backward replacement of these subgames by their minimax
values preserves the root value. No identification of distinct physical
histories is needed.
\end{theorem}
\begin{proof}
At $c$ both previously selected command blocks have finished. Their executed
contents are public, and no private unexecuted command crosses the cut.
All histories belonging to either player's information set at or after this
prefix remain inside the corresponding subtree. The subtree is therefore a
proper subgame. Within it, the maximizing player can guarantee its value
and the minimizing player can enforce the same upper bound. Replacing its
continuation by that scalar gives matching guarantees in the preceding
game. Induction over common boundaries, in reverse chronological order,
proves equality at the root.
\end{proof}
An update of Blue alone is not such a boundary while Red still has an
unfinished command block. This restriction separates the result from an
incorrect full-information Bellman recursion at every physical epoch.

\begin{corollary}[Propagation of subgame error]\label{cor:error}
If the uniform value error at replacement level $l$ is at most $e_l$, the
root error is at most $\sum_l e_l$.
\end{corollary}
\begin{proof}
A uniform perturbation of at most $e$ changes the expected payoff of every
fixed strategy pair by at most $e$. Taking the minimum and maximum preserves
the bound. Induct over the successive replacements.
\end{proof}

\subsection{Exact pricing against a fixed causal opponent}
For a fixed opponent policy $q$, define $C_q(j\mid i)$ as its realization
probability of path $j$ when the own path $i$ is externally fixed. This is a
counterfactual conditional law; it does not assume that the opponent
commits. For committed $q$, it reduces to $q_j$. For a mixture of causal
pure opponent maps $r_m(i)$ with weights $w_m$,
\begin{equation}
 C_q(j\mid i)=\sum_m w_m\mathbf1\{r_m(i)=j\}.
 \label{eq:counterfactual}
\end{equation}
Every positive mixture weight is retained. Initialize
$H_T(i,j)=A(i,j)C_q(j\mid i)$. At consecutive own block boundaries $t<e$,
apply the recurrence
\begin{equation}
 H_t(i_t,j_t)=\max_a\sum_z H_e(i_t a,j_t z),
 \label{eq:pricing}
\end{equation}
where $a$ is the complete own block and $z$ ranges over the opponent's
executed actions during that block. The sum precedes the maximum.

\begin{theorem}[Response over the complete permitted policy class]\label{thm:pricing}
Recurrence \eqref{eq:pricing} returns the maximum payoff against $q$ over
all policies permitted by the own calendar. Its stored maximizing blocks
define a feasible pure best response.
\end{theorem}
\begin{proof}
Causality implies that the counterfactual opponent probability up to time
$t$ depends only on the executed prefix, hence is independent of the
yet-unchosen own block. For each block, all unobserved opponent continuations
must first be summed. At the next own update the continuation can be chosen
separately for each newly observed prefix; its optimum has already been
computed by the induction hypothesis. The resulting sum followed by a
maximum is precisely the feasible decision at $t$. Backward induction proves
root optimality. Prefix probabilities remain unnormalized, so no division
by zero is required. The recorded block at each information set uses only
its observed prefixes, which proves feasibility, including arbitrary legal
choices at zero-weight histories.
\end{proof}
The recurrence takes $O(b^{2T})$ arithmetic and memory for $b\ge2$, up to
a geometric-series factor. Constructing the physical matrix and generating
multiple response policies are additional costs. The result specializes
established extensive-form best-response ideas; we do not claim a new
general equilibrium-computation framework \citep{bosansky2014,mcaleer2021}.

\subsection{Transferable upper bounds and feasible deletion policies}
\begin{proposition}[Opponent-policy reuse]\label{prop:reuse}
For fixed opponent calendar $R$, a causal opponent policy remains feasible
when only the own calendar changes. For any finite pool $Q_0\subseteq\QQ_R$,
\begin{equation}
 V(S,R)\le\min_{q\in Q_0}\BR_S(q).
 \label{eq:poolbound}
\end{equation}
\end{proposition}
\begin{proof}
The opponent's update times, observed executed prefixes and legal action
menu remain unchanged. Its old policy can ignore the new public calendar
and use the same map of observations to commands. Its feasibility never
depends on observing an unexecuted suffix. Fixing any such policy in the
minimax expression gives an upper bound, and the minimum of valid upper
bounds is valid.
\end{proof}

To remove updates, we construct a legal coarse policy rather than treating
an infeasible fine-calendar strategy as if it remained admissible. Express
a full-update source strategy as conditional action kernels. At a coarse
update, retain the actual own and opponent prefixes, complete the opponent
future by a fixed straight-action suffix, and multiply the source kernels
along every possible own block. These products define its block distribution.
When a source information set has zero incoming probability, use the uniform
legal action distribution. This is an explicit new behavioral policy; no
claim is made that it preserves the source's outcome distribution after
updates are deleted.

\begin{theorem}[Constructive deletion certificate]\label{thm:deletion}
Under state-independent legal action menus, the construction gives a feasible
policy $\tilde p$ for any coarse calendar $S'\subseteq\NN$. If $U_F$ bounds
$V_F$ from above and
$L'=\min_{q\in\QQ_R}\EE_{\tilde p,q}A$, then
\begin{equation}
 0\le V_F-V(S',R)\le U_F-L'.
 \label{eq:loss}
\end{equation}
\end{theorem}
\begin{proof}
Each source kernel is normalized. Summing the products over all own blocks
therefore gives one by iterated normalization. The construction uses only
information observed at the coarse update and a fixed hypothetical suffix.
Every block is legal because action menus do not depend on the shadow
trajectory. Thus $\tilde p\in\PP_{S'}$ and $L'\le V(S',R)$. Combine this
inequality with $V_F\le U_F$ and Proposition~\ref{prop:inclusion}.
\end{proof}
The certificate can be loose; its usefulness must be measured. It is a
full-response bound on an explicit policy, not a generic smoothness bound
on the directional field. Existing abstraction theory also distinguishes
solution error from the size of an abstraction \citep{kroer2018}.

\subsection{A calendar bound from interval losses}
The fully replanning opponent admits an additional public-boundary construction.
Let $v_t(h)$ be the full-update continuation value at executed prefix $h$.
For $t<e$, let $b_{t,e}(h)$ be the value when Blue commits over $[t,e)$,
Red remains fully flexible, and the terminal payoff at $e$ is $v_e$.
Define
\begin{equation}
 d(t,e)=\max_{h\text{ at }t}\{v_t(h)-b_{t,e}(h)\}.
 \label{eq:intervalcost}
\end{equation}
Policy inclusion gives $d(t,e)\ge0$, and $d(t,t+1)=0$.

\begin{theorem}[Calendar loss from public intervals]\label{thm:interval}
For opponent $F$ and boundaries $0=t_0<t_1<\cdots<t_m=T$,
\begin{equation}
 0\le V(F,F)-V(\{t_1,\ldots,t_{m-1}\},F)
 \le\sum_{j=0}^{m-1}d(t_j,t_{j+1}).
 \label{eq:intervalloss}
\end{equation}
Minimizing the right side over paths with at most $K+1$ edges from $0$ to $T$
gives a computable upper bound on the loss of its selected calendar.
\end{theorem}
\begin{proof}
At each selected boundary use a maximizing block policy of the game defining
$b_{t,e}(h)$. Suppose the remaining stitched policy guarantees at least
$v_e(h')-D$ at every successor prefix $h'$. Replacing $v_e$ by that
continuation lowers the guarantee by at most $D$. The selected block thus
guarantees $b_{t,e}(h)-D\ge v_t(h)-d(t,e)-D$. Backward induction from zero
terminal loss proves the sum. Every stitched decision uses an observed
public prefix; Red has no private unexecuted block across these cuts.
Finally each forward path is exactly one feasible public calendar, so
shortest-path minimization preserves the bound.
\end{proof}
With edge costs available, all budget bounds take $O(T^3)$ dynamic-programming
work. Computing the full value tree and local interval games is preprocessing
and can be expensive. Maxima over all histories can be conservative.
The bound certifies sufficiency of a budget, not minimality for the true
performance target. The construction is restricted to opponent $F$;
an unexecuted committed opponent suffix prevents the same public-boundary
argument for $C$. The result specializes a minimax perturbation argument
to calendar selection rather than claiming a new general abstraction theorem.

\subsection{Correct use of symmetry}
Exchange-equivariant dynamics and antisymmetric payoff imply
$V(\mathsf E s)=-V(s)$ after exchanging policy classes. At states fixed by
the exchange, possibly combined with a verified spatial isometry, this
implies zero value when policy classes are also exchanged into themselves.
It does not imply that the finite-horizon value equals one-stage payoff at
every state. For example, absorbing states $s=\pm1$ with stage payoff $s$
have two-stage value $2s$ while satisfying exchange oddness. Mixed minimax,
rather than pure max-min, is also essential for the zero value of a
skew-symmetric matrix. These distinctions govern the numerical symmetry
checks and prevent an invalid degeneracy argument from entering the main
line of reasoning.
